\documentclass[10pt,a4paper]{amsart}
\usepackage[margin=19mm]{geometry}
\usepackage{amsmath,amssymb,mathtools}
\usepackage[scr=rsfs,cal=euler]{mathalfa}
\usepackage{xfrac}
\usepackage[unicode,hidelinks,bookmarks=false]{hyperref}
\newcommand{\ie}{i.e.\ }
\newcommand{\cf}{cf.\ }
\newcommand{\resp}{respectively\ }
\newcommand{\Subsection}{\subsection}
\providecommand{\nobreakdash}{\nobreak\hbox{-}}
\setlength{\emergencystretch}{2em}
\multlinegap0pt
\usepackage{bm}
\usepackage{bbm}
\usepackage{dsfont}
\usepackage{xspace}
\usepackage{cancel}
\usepackage{mathtools}
\usepackage{amsthm}
\makeatletter
\@ifundefined{theorem}{\newtheorem{theorem}{Theorem}}{}
\@ifundefined{lemma}{\newtheorem{lemma}[theorem]{Lemma}}{}
\makeatother
\newcommand{\Ai}{\operatorname{Ai}}
\newcommand{\Bi}{\operatorname{Bi}}
\title[The Tricomi equation in the hyperbolic half plane under addi]{The Tricomi equation in the hyperbolic half plane under additive space-time Gaussian White Noise perturbation}
\author{Enrico Bernardi, Alberto Lanconelli}
\date{Source: arXiv:2606.04950v1 (CC BY 4.0)}
\begin{document}
\maketitle

\noindent\emph{Source and licence.} The Tricomi equation in the hyperbolic half plane under additive space-time Gaussian White Noise perturbation. Version arXiv:2606.04950v1 (CC BY 4.0), distributed under the Creative Commons Attribution 4.0 International licence, as recorded in the arXiv licence metadata for this exact version and shown on the abstract page https://arxiv.org/abs/2606.04950v1. Full rights notice: licenses/arxiv-2606.04950-CC-BY-4.0.txt.\par\smallskip

\noindent\textit{Editorial scope.} Counted unit: the Lemma whose source label is \texttt{as} in the article (source0.tex lines 493--518, source label \texttt{as}), delivered together with the article's own complete original proof of it (source0.tex lines 520--845, one \texttt{proof} environment of 326 lines). No mathematics is written, repaired or re-derived: statement and proof are the article's own text line for line. Required context retained uncounted: none. Every definition, display and label of those lines is retained unchanged. Omitted from this package and not redistributed: 1--492, 519--519, 846--1111. Nothing inside a retained excerpt is omitted or altered: every retained body is delivered exactly as the article's lines are written, so no omission receipt of any kind accompanies this package. Citations: the retained excerpts carry no citation command at all, so no bibliography entry of the article is transcribed and no reference is redistributed. Reference adaptation: none was needed and none was made; every reference inside the delivered unit resolves inside it, so no unresolved reference or citation is produced. Printed slips kept literal: no printed word, symbol, hypothesis or number of the counted statement or of its proof was altered. Layout and class: the article's own class is \texttt{article} with packages including amsmath, amssymb, amsthm, tikz, color, appendix, graphicx, fancyhdr, enumitem, bbm, parskip, float, chngpage, calc; the wrapper is the lane's pinned \texttt{amsart} editorial wrapper at 10pt on A4, which restates the theorem-like environments the retained text needs and adds the article's own macro definitions that the retained text needs (\texttt{\textbackslash Ai}, \texttt{\textbackslash Bi}), with extra wrapper packages (bm, bbm, dsfont, xspace, cancel, mathtools). The delivered numbering is the unit's own. Assets: no retained span contains a figure environment, a graphics-inclusion command, a PDF-page inclusion, a table or a TikZ picture; no image file is copied into this package, the package declares no asset dependency and no image file is redistributed with it. Licence: version arXiv:2606.04950v1 of this article is distributed under the Creative Commons Attribution 4.0 International licence, as recorded in the arXiv licence metadata for this exact version and shown on the abstract page https://arxiv.org/abs/2606.04950v1. A full rights notice is delivered at licenses/arxiv-2606.04950-CC-BY-4.0.txt. Characters: every retained excerpt is pure ASCII, so the delivered engine, XeLaTeX with its default Latin Modern text font, reports no missing character.
\begin{lemma}\label{as}
Fix \(t>0\), and set
\[
  K_t(\rho):=\int_0^t |C(t,s,\rho)|^2\,ds ,
\]
where
\[
C(t,s,\rho)
=
-\pi\rho^{-2/3}
\Big[
 \Bi(-\rho^{2/3}t)\Ai(-\rho^{2/3}s)
 -
 \Ai(-\rho^{2/3}t)\Bi(-\rho^{2/3}s)
\Big].
\]
Then
\[
  K_t(\rho)=\rho^{-2}+o(\rho^{-2})
  \qquad \text{as } \rho\to+\infty.
\]
Equivalently,
\[
  \rho^2K_t(\rho)\longrightarrow 1 .
\]
\end{lemma}

\begin{proof}
Fix \(t>0\). We shall prove that
\[
  \rho^2\int_0^t |C(t,s,\rho)|^2\,ds \longrightarrow 1 .
\]
We use the standard Airy asymptotics on the negative real axis. Namely, if
\[
  \theta(x):=\frac23 x^{3/2}+\frac{\pi}{4},
  \qquad x>0,
\]
then, as \(x\to+\infty\),
\[
  \Ai(-x)
  =
  \pi^{-1/2}x^{-1/4}
  \bigl(\sin\theta(x)+r_A(x)\bigr),
\]
and
\[
  \Bi(-x)
  =
  \pi^{-1/2}x^{-1/4}
  \bigl(\cos\theta(x)+r_B(x)\bigr),
\]
where
\[
  r_A(x)\to0,\qquad r_B(x)\to0
  \qquad \text{as } x\to+\infty.
\]
Equivalently, if
\[
  \varepsilon_L
  :=
  \sup_{x\ge L}\bigl(|r_A(x)|+|r_B(x)|\bigr),
\]
then
\[
  \varepsilon_L\to0
  \qquad \text{as } L\to+\infty.
\]
Let \(L>1\) be fixed. For \(\rho\) large enough we have
\[
  L\rho^{-2/3}<t.
\]
We split
\[
  \rho^2K_t(\rho)
  =
  \rho^2\int_0^{L\rho^{-2/3}} |C(t,s,\rho)|^2\,ds
  +
  \rho^2\int_{L\rho^{-2/3}}^t |C(t,s,\rho)|^2\,ds
  =: I_{\mathrm{small}}(\rho,L)+I_{\mathrm{bulk}}(\rho,L).
\]
We first show that the small interval gives no contribution. On
\[
  0\le s\le L\rho^{-2/3},
\]
the quantities \(\rho^{2/3}s\) remain in the compact interval \([0,L]\). Hence
\(\Ai(-\rho^{2/3}s)\) and \(\Bi(-\rho^{2/3}s)\) are bounded by a constant depending only on \(L\). On the other hand,
\[
  \Ai(-\rho^{2/3}t)=O((\rho^{2/3}t)^{-1/4})
  =O(\rho^{-1/6}),
\]
and similarly
\[
  \Bi(-\rho^{2/3}t)=O(\rho^{-1/6}).
\]
Therefore
\[
\begin{aligned}
 &\Big|
 \Bi(-\rho^{2/3}t)\Ai(-\rho^{2/3}s)
 -
 \Ai(-\rho^{2/3}t)\Bi(-\rho^{2/3}s)
 \Big|
 \le C_L\rho^{-1/6}.
\end{aligned}
\]
Consequently,
\[
  |C(t,s,\rho)|^2
  \le C_L \rho^{-4/3}\rho^{-1/3}
  =
  C_L\rho^{-5/3}.
\]
Thus
\[
  I_{\mathrm{small}}(\rho,L)
  \le
  \rho^2
  \int_0^{L\rho^{-2/3}} C_L\rho^{-5/3}\,ds
  =
  C_L L \rho^{-1/3}
  \longrightarrow 0
\]
as \(\rho\to+\infty\), for every fixed \(L\).\\
We now study \(I_{\mathrm{bulk}}(\rho,L)\). Put
\[
  x_t:=\rho^{2/3}t,\qquad x_s:=\rho^{2/3}s .
\]
On the interval \(s\in[L\rho^{-2/3},t]\), both \(x_s\) and \(x_t\) are at least \(L\), for \(\rho\) sufficiently large. Therefore the Airy asymptotics above apply uniformly. Define
\[
\begin{aligned}
D(t,s,\rho)
&:=
\Bi(-x_t)\Ai(-x_s)
-
\Ai(-x_t)\Bi(-x_s).
\end{aligned}
\]
Using the asymptotic formulae, we get
\[
\begin{aligned}
D(t,s,\rho)
&=
\pi^{-1}(x_tx_s)^{-1/4}
\Big[
 \cos\theta(x_t)\sin\theta(x_s)
 -
 \sin\theta(x_t)\cos\theta(x_s)
 +
 R_L(t,s,\rho)
\Big],
\end{aligned}
\]
where, uniformly for \(s\in[L\rho^{-2/3},t]\),
\[
  |R_L(t,s,\rho)|\le C\varepsilon_L .
\]
Since
\[
  \cos\theta(x_t)\sin\theta(x_s)
  -
  \sin\theta(x_t)\cos\theta(x_s)
  =
  \sin(\theta(x_s)-\theta(x_t)),
\]
and since
\[
  \theta(x_t)-\theta(x_s)
  =
  \frac23 \rho(t^{3/2}-s^{3/2}),
\]
we obtain
\[
D(t,s,\rho)
=
\pi^{-1}\rho^{-1/3}(ts)^{-1/4}
\Big[
 -\sin\Big(\frac23\rho(t^{3/2}-s^{3/2})\Big)
 +
 R_L(t,s,\rho)
\Big].
\]
Multiplying by the prefactor in \(C\), namely \(C=-\pi\rho^{-2/3}D\), gives
\[
C(t,s,\rho)
=
\rho^{-1}(ts)^{-1/4}
\Big[
 \sin\Big(\frac23\rho(t^{3/2}-s^{3/2})\Big)
 +
 \widetilde R_L(t,s,\rho)
\Big],
\]
with
\[
  |\widetilde R_L(t,s,\rho)|\le C\varepsilon_L
\]
uniformly on the same interval. Therefore
\[
\begin{aligned}
I_{\mathrm{bulk}}(\rho,L)
&=
\int_{L\rho^{-2/3}}^t
(ts)^{-1/2}
\Big|
 \sin\Big(\frac23\rho(t^{3/2}-s^{3/2})\Big)
 +
 \widetilde R_L(t,s,\rho)
\Big|^2\,ds .
\end{aligned}
\]
It follows that
\[
\begin{aligned}
I_{\mathrm{bulk}}(\rho,L)
&=
J(\rho,L)+E(\rho,L),
\end{aligned}
\]
where
\[
  J(\rho,L)
  :=
  \int_{L\rho^{-2/3}}^t
  (ts)^{-1/2}
  \sin^2\Big(\frac23\rho(t^{3/2}-s^{3/2})\Big)\,ds,
\]
and
\[
  |E(\rho,L)|
  \le
  C\varepsilon_L
  \int_{L\rho^{-2/3}}^t (ts)^{-1/2}\,ds .
\]
But
\[
  \int_{L\rho^{-2/3}}^t (ts)^{-1/2}\,ds
  =
  t^{-1/2}
  \int_{L\rho^{-2/3}}^t s^{-1/2}\,ds
  \le 2.
\]
Hence
\[
  |E(\rho,L)|\le C\varepsilon_L .
\]
It remains to compute the limit of \(J(\rho,L)\). Using
\[
  \sin^2 y=\frac12(1-\cos(2y)),
\]
we write
\[
\begin{aligned}
J(\rho,L)
&=
\frac12 t^{-1/2}
\int_{L\rho^{-2/3}}^t s^{-1/2}\,ds        \\
&\quad
-
\frac12 t^{-1/2}
\int_{L\rho^{-2/3}}^t
s^{-1/2}
\cos\Big(\frac43\rho(t^{3/2}-s^{3/2})\Big)\,ds .
\end{aligned}
\]
The non-oscillatory part is
\[
  \frac12 t^{-1/2}
  \int_{L\rho^{-2/3}}^t s^{-1/2}\,ds
  =
  t^{-1/2}
  \left(\sqrt t-\sqrt L\,\rho^{-1/3}\right)
  \longrightarrow 1 .
\]
We claim that the oscillatory part tends to zero. Set
\[
  u=\frac23(t^{3/2}-s^{3/2}).
\]
Then
\[
  du=-s^{1/2}\,ds,
  \qquad
  s=\left(t^{3/2}-\frac32u\right)^{2/3},
\]
and therefore
\[
  s^{-1/2}\,ds
  =
  -s^{-1}\,du
  =
  -\left(t^{3/2}-\frac32u\right)^{-2/3}\,du .
\]
Thus
\[
\begin{aligned}
&\int_{L\rho^{-2/3}}^t
s^{-1/2}
\cos\Big(\frac43\rho(t^{3/2}-s^{3/2})\Big)\,ds      \\
&\qquad =
\int_0^{\frac23(t^{3/2}-L^{3/2}\rho^{-1})}
\left(t^{3/2}-\frac32u\right)^{-2/3}
\cos(2\rho u)\,du .
\end{aligned}
\]
The function
\[
  g(u):=
  \left(t^{3/2}-\frac32u\right)^{-2/3}
\]
belongs to
\[
  L^1\left(0,\frac23t^{3/2}\right),
\]
because its singularity at the endpoint is of order \(2/3<1\). Hence, by the Riemann--Lebesgue lemma,
\[
  \int_0^{\frac23t^{3/2}}
  g(u)\cos(2\rho u)\,du
  \longrightarrow 0 .
\]
Moreover,
\[
\begin{aligned}
&\left|
\int_{\frac23(t^{3/2}-L^{3/2}\rho^{-1})}^{\frac23t^{3/2}}
g(u)\cos(2\rho u)\,du
\right|        \\
&\qquad \le
\int_{\frac23(t^{3/2}-L^{3/2}\rho^{-1})}^{\frac23t^{3/2}}
g(u)\,du
\longrightarrow 0 .
\end{aligned}
\]
Therefore the oscillatory part tends to zero, and so
\[
  J(\rho,L)\longrightarrow 1
  \qquad \text{as } \rho\to+\infty .
\]
We have shown that, for every fixed \(L>1\),
\[
  \limsup_{\rho\to+\infty}
  \bigl|\rho^2K_t(\rho)-1\bigr|
  \le C\varepsilon_L .
\]
Finally, since \(\varepsilon_L\to0\) as \(L\to+\infty\), we conclude that
\[
  \rho^2K_t(\rho)\longrightarrow 1 .
\]
Equivalently,
\[
  K_t(\rho)=\rho^{-2}+o(\rho^{-2})
  \qquad \text{as } \rho\to+\infty .
\]
The proof is complete.
\end{proof}

\end{document}
